\section{基于运行生产资料的汽轮机系统故障知识图谱构建}
\label{sec:knowledge_graph_v2}

\subsection{引言}

上一章基于正常运行数据建立了燃煤机组汽轮机系统设备级状态监测方法，其作用是在给定运行边界和操作状态下获得设备的正常状态参考。然而，状态监测模型主要解决“运行状态是否偏离健康参考”的问题，并不能仅依靠参数偏差直接给出设备归属、故障机理和传播路径。对于由多级汽缸、再热系统、抽汽回热系统、给水系统和凝汽器冷端共同组成的汽轮机热力系统，不同设备之间存在明确的蒸汽流动、能量交换和控制关系，一个局部故障可能首先改变设备内部状态，随后沿热力接口影响相邻系统。若将所有运行参数视为相互独立的特征输入分类器，虽然能够学习数据中的统计相关性，却难以显式表达这些系统结构和工程知识。

知识图谱（Knowledge Graph, KG）以实体、关系和属性组织异构知识，能够将分散于系统结构、运行规程、设备资料、故障案例和学术文献中的知识表示为机器可计算的图结构\cite{Hogan2021KnowledgeGraphs}。在工业故障诊断场景中，知识图谱并不替代实时运行数据，而是为诊断模型提供“参数属于哪个设备”“设备之间如何连接”“某类故障会引起哪些异常”“哪些监测参数能够反映相应异常”等结构先验。Han等将过程知识与图神经网络结合，用于约束复杂工业过程中的故障传播和根因分析，说明领域知识能够减少纯数据驱动图结构中的不合理关系\cite{Han2022KnowledgeEnhancedGNN}。针对汽轮机系统，Zhang和Hu根据设备机理构建过程图与故障传播图，并利用图神经网络推断汽封缺陷、蒸汽泄漏和汽耗增加之间的传播路径，验证了图结构表示汽轮机故障传播关系的可行性\cite{ZhangHu2021FaultPropagationGNN}。

基于上述认识，本章借鉴课题组已有运行生产资料知识图谱构建思路\cite{YangHaodong2026CFB}，结合汽轮机系统特点重新定义知识范围、实体模式和关系模式。首先从汽轮机主通流、再热、抽汽回热、给水及冷端系统分析典型热力故障及其参数响应；随后组织系统拓扑、DCS测点、设备资料和公开同行评议文献中的故障知识，构建“设备--参数--异常--故障--处置”五类实体组成的领域语义图；对于后续大规模运行日志、缺陷记录和检修文本，则设计“大语言模型候选抽取--领域规则校验--语义对齐--证据评价”的知识扩展流程。最后建立领域参数概念与实际DCS测点之间的映射，为下一章KG-GCN模型提供静态图拓扑基础。

需要说明的是，本章当前构建的核心知识图谱以实际系统拓扑、DCS参数定义和逐篇原始文献证据核验为基础，属于经证据审核后的核心语义层；大语言模型和Embedding方法主要用于后续大规模文本知识扩展与实体归一化。本文不将未经证据核验的模型生成关系直接写入核心图谱，也不将跨机型文献中的故障关系等同于本文机组已发生的现场故障。

\subsection{燃煤机组汽轮机系统参数及故障传播特性}
\label{subsec:turbine_fault_characteristics_v2}

\subsubsection{系统拓扑与故障传播路径}

燃煤机组汽轮机系统具有典型的层次化和网络化结构。主通流沿“主蒸汽--VHP--一次再热RH1--HP--二次再热RH2--IP--LP--凝汽器”的方向传递，抽汽从不同压力等级汽缸分流至高、低压加热器和除氧器，凝结水经低压加热、除氧、给水泵升压和高压加热后重新进入锅炉，凝汽器则通过循环水系统形成冷端边界。上述主通流、抽汽回热、给水和冷端回路构成了汽轮机系统中最主要的物质和能量传播路径。

故障对运行状态的影响通常具有局部产生、沿接口扩散的特点。首先，故障会改变故障设备内部原有的热力状态，例如汽封泄漏会改变有效蒸汽流量和汽缸效率，叶片侵蚀或沉积会改变通流面积和级组压力分布；随后，这些变化可能通过排汽、抽汽或冷端接口作用于相邻设备，例如VHP排汽状态影响一次再热入口，抽汽热力状态影响对应加热器换热过程，凝汽器背压则反向影响LP末级膨胀。Zhang和Hu指出，大型工业系统中的故障可沿物质流或信息流路径传播，在汽轮机泄漏算例中构建了由汽封缺陷到蒸汽泄漏再到汽耗增加的分层传播图\cite{ZhangHu2021FaultPropagationGNN}。因此，汽轮机知识图谱的关系边不应仅由统计相关性决定，而应优先来自实际热力拓扑和经过证据支持的故障机理。

按照知识在诊断过程中的作用，可将本章关系分为三类。第一类为\textbf{结构关系}，用于描述设备归属、蒸汽流向、抽汽连接和冷端连接；第二类为\textbf{故障机理关系}，用于描述“故障--异常--参数”之间的因果、影响和表现链；第三类为\textbf{诊断处置关系}，用于描述故障与监测参数、运行处置及检修动作之间的联系。三类关系共同形成从系统结构、状态变化到故障诊断的知识链路。

\subsubsection{典型热力故障及参数响应特征}

汽轮机故障既包括轴系振动、转子不平衡等机械故障，也包括通流、汽封、阀门、回热和冷端异常。由于本文前述状态监测模型主要利用DCS热力变量，后续故障诊断也以通流和热力状态为重点，因此本章核心知识图谱不将机械振动故障作为主要研究对象，而优先选择能够由压力、温度、流量、阀位和效率等参数反映的故障类型。

（1）\textbf{调节阀死区及配汽异常。} 高压调节阀负责调节进入汽轮机的蒸汽流量并参与机组功率控制。阀杆、阀体和执行机构长期频繁动作可能形成磨损、卡涩或死区。Zhang和Hu针对汽轮机高压调节阀建立图模型，指出死区表现为控制指令变化而实际阀位在一定范围内不发生响应，并可能引起高压侧和系统振荡\cite{ZhangHu2021ControlValveGraph}。因此可形成“调节阀死区--导致--阀位对指令不敏感--表现为--阀位反馈异常”的故障链，并将主蒸汽流量及入口压力波动作为关联监测信息。

（2）\textbf{汽封磨损和蒸汽泄漏。} 级间汽封和轴端汽封用于限制不同压力区域之间的蒸汽泄漏。汽封齿损伤、间隙增大或部件磨损会增加泄漏量。汽轮机故障传播研究表明，HP、IP和LP的级间或轴端汽封缺陷可导致相应部位蒸汽泄漏，并进一步表现为维持相同输出功率所需汽耗增加；对于LP轴端密封异常，还可能涉及空气漏入低压系统并影响凝汽器真空\cite{ZhangHu2021FaultPropagationGNN}。针对600~MW超临界汽轮机HP--IP中间汽封泄漏的热力研究进一步指出，中间汽封泄漏是影响IP效率的重要因素，热再温度和IP排汽温度是泄漏量估计中的关键测量参数\cite{Xu2011HPIPLeakage}。

（3）\textbf{通流部件侵蚀与沉积。} 长期运行过程中，汽轮机通流部件可能出现叶片表面粗糙度增加、固体颗粒侵蚀（Solid Particle Erosion, SPE）、汽封磨损和叶片表面沉积等退化。Kubiak等基于多台汽轮机现场诊断及检修验证指出，这些现象属于常见的通流性能退化形式；其中SPE可能造成首级喷嘴有效面积增大并改变关键压力分布\cite{Kubiak2002TurbineDegradation}。叶片表面沉积会减小有效通流面积并导致出力和效率下降，级前后压力可用于反映通流面积变化\cite{Kubiak2002ScaleDeposit}。需要强调的是，后者研究对象包含不同类型汽轮机，本文将其作为通用通流机理知识使用，只有与本文机组设备结构和可观测参数兼容的关系才进入领域语义层，而不据此声称本文机组已经发生相同故障。

（4）\textbf{凝汽器污垢及冷端性能劣化。} 凝汽器水侧污垢会增加传热热阻并降低换热能力。Li等针对600~MW超临界燃煤机组构建凝汽器在线污垢监测模型，指出水侧沉积会降低凝汽器热有效性，并可利用污垢热阻变化实现在线监测\cite{Li2016CondenserFouling}。Medica-Viola等进一步表明，凝汽压力受到汽轮机负荷、循环水流量、循环水入口温度和换热面污垢等因素共同影响；凝汽压力升高会降低汽轮机可利用膨胀焓降并影响效率和出力\cite{Medica2018CondenserCondition}。因此，知识图谱中既需要表达“凝汽器污垢--换热性能下降--背压升高”的故障链，也需要保留循环水温度和流量等外部边界，避免将正常环境变化误解释为设备故障。

（5）\textbf{给水加热器内部泄漏及回热性能下降。} 给水加热器内部管束、管板或隔板泄漏会改变正常的抽汽--给水换热过程。Barszcz和Czop针对燃煤机组给水加热器提出用于模型诊断的动态模型，并指出换热功率等过程指标可用于跟踪包括管道裂纹内部泄漏在内的慢性性能变化\cite{Barszcz2011FeedwaterHeater}。Heo和Lee研究了闭式给水加热器内部泄漏的定位及泄漏量诊断，表明内部泄漏会影响回热循环效率并与相邻汽轮机和凝汽器状态形成耦合\cite{Heo2012FWHLeakage}。该研究对象的机组类型与本文不同，因此相应关系仅作为换热设备通用机理知识，最终需通过本文机组的设备形式、运行规程和DCS测点进行兼容性审核。

表~\ref{tab:turbine_fault_knowledge_examples}汇总了当前核心知识图谱中具有较强文献证据的典型故障链。可以看到，不同故障虽然发生部位不同，但均可归纳为“故障实体--异常实体--监测参数--影响设备”的多层关系。这一共性直接决定了后续知识图谱需要同时设置设备、参数、异常和故障等不同语义实体，而不能仅构造参数之间的单一相关网络。

\begin{table}[htbp]
\centering
\caption{典型汽轮机热力故障及知识关系示例}
\label{tab:turbine_fault_knowledge_examples}
\begin{tabular}{p{2.7cm}p{3.7cm}p{4.2cm}p{3.3cm}}
\toprule
故障/退化 & 主要异常或机理 & 关键监测信息 & 主要影响对象 \\
\midrule
高压调节阀死区 & 阀位对指令不敏感、系统振荡 & 阀位反馈、主蒸汽流量、入口压力波动 & VHP及调节系统 \\
汽封缺陷/磨损 & 蒸汽泄漏、汽耗增加 & 热再温度、IP排汽温度、相关压力/流量 & HP、IP、LP \\
固体颗粒侵蚀 & 喷嘴面积变化、压力分布变化 & 级前后压力、热再/连通管压力 & 汽轮机通流级组 \\
叶片表面沉积 & 通流面积减小、效率和出力下降 & 级前后压力、出力、效率 & 汽轮机通流级组 \\
凝汽器水侧污垢 & 换热能力下降、背压升高 & 污垢热阻、背压、循环水温度/流量 & 凝汽器、LP \\
给水加热器内漏 & 回热性能下降 & 给水出口温度、抽汽侧压力、换热指标 & 回热系统、机组效率 \\
\bottomrule
\end{tabular}
\end{table}

\subsection{基于运行生产资料的知识图谱构建方法}
\label{subsec:kg_construction_method_v2}

\subsubsection{知识来源与证据分级}

知识图谱的可靠性首先取决于知识来源。杨昊东针对CFB锅炉运行生产资料的研究利用设计资料、运行日志和缺陷记录构造故障知识，并通过大语言模型、规则校验和语义对齐完成结构化处理\cite{YangHaodong2026CFB}。面向本文汽轮机系统，考虑当前现场历史故障资料尚不完整，本阶段采用“核心语义层先行、运行文本持续扩展”的构建方式。不同资料的主要作用见表~\ref{tab:kg_evidence_sources}。

\begin{table}[htbp]
\centering
\caption{汽轮机故障知识来源及主要用途}
\label{tab:kg_evidence_sources}
\begin{tabular}{p{3.0cm}p{5.2cm}p{5.2cm}}
\toprule
资料类型 & 主要提取知识 & 本文使用方式 \\
\midrule
系统图、DCS画面与点表 & 设备、参数、归属、主通流、抽汽和冷端连接 & 作为本机组结构关系的最高优先级依据 \\
运行规程、设备说明书 & 设备功能、正常边界、报警含义、运行处置 & 后续随现场资料补充扩展核心图谱 \\
缺陷/检修/运行日志 & 故障现象、故障原因、处置动作、共现关系 & 采用LLM辅助抽取并要求原文证据回溯 \\
同行评议原始论文 & 通用故障机理、参数响应、传播关系 & 兼容本文机组结构后进入领域语义层 \\
既有专家知识库 & 别名、异常描述、处置经验 & 作为辅助证据，不单独形成高置信因果边 \\
\bottomrule
\end{tabular}
\end{table}

本章当前核心语义图谱主要由三部分知识构成：一是根据本机组系统拓扑和DCS资料确定的设备--参数和设备--设备结构关系；二是从原始同行评议论文中逐条核验得到的故障--异常--参数关系；三是依据标准术语表进行的实体归一化。对于来源于其他容量等级、核电二回路或地热汽轮机等场景的文献关系，仅提取不依赖特定机型的通用热力机理，并再次检查本文机组是否存在对应设备和可观测参数。因而，“文献支持”在本章表示该知识关系具有外部机理证据，而不表示该故障已在本文机组现场数据中得到验证。

对于后续大量运行日志、缺陷记录和检修文本，采用自动化知识扩展流程。大语言模型负责从非结构化文本中生成候选实体和三元组，规则层负责检查实体类型、关系方向、设备拓扑和测点存在性，Embedding用于别名归一化与语义对齐，最终保留原始证据位置供人工或专家复核。该流程的目标是降低大规模文本人工整理成本，而不是利用语言模型生成缺少依据的新故障知识。

\begin{figure}[htbp]
\centering
\fbox{\parbox[c][4.2cm][c]{0.90\textwidth}{\centering
图3-1占位：多源运行生产资料\\
系统拓扑 / DCS点表 / 运行规程 / 故障与检修记录 / 同行评议文献\\
$\downarrow$ LLM候选抽取 $\rightarrow$ 规则校验 $\rightarrow$ 实体标准化 $\rightarrow$ 证据审核 $\rightarrow$ 汽轮机故障KG}}
\caption{汽轮机故障知识抽取与图谱构建流程}
\label{fig:kg_construction_flow}
\end{figure}

\subsubsection{实体类型与关系模式定义}

本文将汽轮机故障知识图谱表示为有向多关系图
\begin{equation}
\mathcal{G}_{sem}=(\mathcal{V},\mathcal{E},\mathcal{R}),
\label{eq:kg_definition_v2}
\end{equation}
其中，$\mathcal{V}$为实体集合，$\mathcal{R}$为关系类型集合，每条有向关系边表示为$e_k=(v_i,r_k,v_j)\in\mathcal{E}$。图在后续GCN计算中是否进行反向边补充或邻接矩阵对称化，将根据第四章具体输入构建方式处理，本章保留原始知识关系方向。

借鉴课题组已有知识图谱实体组织方式\cite{YangHaodong2026CFB}，并结合汽轮机故障诊断需求，实体集合划分为设备、参数、异常、故障和处置五类：
\begin{equation}
\mathcal{V}=\mathcal{V}_{D}\cup\mathcal{V}_{P}\cup\mathcal{V}_{A}\cup\mathcal{V}_{F}\cup\mathcal{V}_{O}.
\end{equation}
设备实体$\mathcal{V}_{D}$包括VHP、RH1、HP、RH2、IP、LP、凝汽器、加热器和阀门等；参数实体$\mathcal{V}_{P}$包括压力、温度、流量、阀位、效率等可测或可计算量；异常实体$\mathcal{V}_{A}$描述蒸汽泄漏、真空下降、通流面积变化和系统振荡等偏离正常状态的现象；故障实体$\mathcal{V}_{F}$描述导致异常的具体故障或性能退化原因；处置实体$\mathcal{V}_{O}$表示运行调整或检修动作。

关系类型主要包括：参数到设备的\textbf{归属}关系；设备之间的\textbf{介质流向}和\textbf{介质/能量连接}；故障到设备的\textbf{发生于}；故障或异常之间的\textbf{导致}和\textbf{影响}；抽象异常到可观测量的\textbf{表现为}；故障到诊断参数的\textbf{关联监测}；以及故障到运行/检修行为的\textbf{处置为}。通过实体--关系schema限定，可以减少同一自然语言关系被多种不一致形式表达的问题，并为后续从语义图谱中提取参数邻接子图提供统一接口。

\begin{figure}[htbp]
\centering
\fbox{\parbox[c][4.0cm][c]{0.90\textwidth}{\centering
图3-2占位：五类实体及关系schema\\
设备 $\leftrightarrow$ 参数；设备 $\rightarrow$ 设备；故障 $\rightarrow$ 异常 $\rightarrow$ 参数；故障 $\rightarrow$ 处置}}
\caption{汽轮机故障知识图谱实体与关系模式示意图}
\label{fig:kg_schema}
\end{figure}

\subsubsection{候选三元组抽取与领域规则校验}

面向后续运行日志、缺陷记录和检修文本，首先利用大语言模型进行候选知识抽取。文本片段$d_k$经提示模板约束后输出候选三元组
\begin{equation}
\tau_k=(h_k,r_k,t_k,s_k),
\end{equation}
其中$h_k$和$t_k$分别为头实体和尾实体，$r_k$为候选关系，$s_k$记录原始证据文本及来源。该阶段输出仅视为候选知识，不直接进入核心图谱。

领域规则层进一步从五个方面进行检查：（1）\textbf{实体类型规则}，例如“发生于”的尾实体必须为设备或部件；（2）\textbf{关系方向规则}，统一“参数--归属--设备”和“上游设备--介质流向--下游设备”等关系方向；（3）\textbf{设备拓扑规则}，传播路径应与实际主通流、抽汽、给水和冷端结构兼容；（4）\textbf{测点存在性规则}，涉及本文机组实时参数时必须能够在标准点表中找到对应测点或明确的领域参数概念；（5）\textbf{证据回溯规则}，故障因果边必须保留原始资料或文献出处。通过规则约束，能够将语言模型开放式抽取结果限制在汽轮机领域允许的语义结构内。

当前107节点、114条三元组的核心语义层并非由上述自动流程一次性生成，而是采用“系统拓扑/DCS概念整理+文献检索+逐条证据核验+schema规则化”的方式构建。自动抽取流程主要用于后续现场文本规模扩大后的持续扩展。这一区分能够保证本章当前报告的图谱规模与实际完成工作一致。

\subsubsection{实体标准化与关系证据评价}

同一设备或参数在不同资料中可能存在多种表述。例如“凝汽器压力”“凝汽器背压”和“排汽背压”在特定语境中具有相近含义，“调节阀反馈”“阀位反馈”以及具体KKS点名也可能指向相同参数概念。若直接保留全部文本名称，会造成大量重复节点并影响DCS映射。因此，本文建立标准词表，并按照“本机组KKS/点表标准名--设备标准名称--文献与运行文本别名”的优先顺序进行实体归一化。

对候选实体$e_i$和标准实体$e_j$，分别计算字符串相似度$S_{str}$和语义相似度$S_{sem}$。语义向量可利用Sentence-BERT等模型获得\cite{Reimers2019SentenceBERT}，综合相似度写为
\begin{equation}
S(e_i,e_j)=\lambda S_{sem}(e_i,e_j)+(1-\lambda)S_{str}(e_i,e_j).
\label{eq:entity_similarity_v2}
\end{equation}
当综合相似度超过阈值、实体类型一致且设备上下文不冲突时，将候选实体映射到标准实体；对于KKS点名、相邻设备同名测点和专业缩写，保留人工核验。当前研究尚未针对大规模现场文本标定最佳$\lambda$和阈值，因此式~\eqref{eq:entity_similarity_v2}作为后续自动扩展的实体对齐框架，而不在现阶段虚构固定参数。

关系证据评价同样采用多源信息。借鉴既有知识图谱构建思路，可将候选关系评分表示为
\begin{equation}
C_k=w_1C_{model,k}+w_2C_{rule,k}+w_3C_{sem,k}+w_4C_{evidence,k},\quad \sum_{i=1}^{4}w_i=1,
\label{eq:relation_confidence_v2}
\end{equation}
其中$C_{model,k}$表示大语言模型候选抽取评分，$C_{rule,k}$表示领域规则满足程度，$C_{sem,k}$表示实体语义一致性，$C_{evidence,k}$表示独立资料、文献或现场记录的证据支持程度。式~\eqref{eq:relation_confidence_v2}描述的是后续自动知识扩展的筛选框架；当前核心语义图谱尚未对全部114条关系使用统一权重进行定量评分，而是按照来源可追溯性将关系记录为直接证据和辅助证据等级，避免将尚未标定的评分模型包装成已完成实验结果。

\subsection{汽轮机系统故障知识图谱构建结果}
\label{subsec:kg_result_v2}

\subsubsection{核心语义知识图谱}

依据上述schema和证据原则，本文首先构建汽轮机系统故障知识图谱的核心语义层。该图谱覆盖主通流、再热、抽汽回热、给水和冷端系统，重点收录能够由热力运行参数反映的通流/性能故障。经实体标准化和关系去重后，当前核心语义层包含107个节点和114条证据三元组，其中设备实体23个、参数实体39个、异常实体24个、故障实体15个、处置实体6个。主要关系数量见表~\ref{tab:kg_relation_stats}。

\begin{table}[htbp]
\centering
\caption{汽轮机故障核心语义图谱实体规模}
\label{tab:kg_entity_stats_v2}
\begin{tabular}{lccccc}
\toprule
实体类型 & 设备 & 参数 & 异常 & 故障 & 处置 \\
\midrule
节点数 & 23 & 39 & 24 & 15 & 6 \\
\bottomrule
\end{tabular}
\end{table}

\begin{table}[htbp]
\centering
\caption{汽轮机故障核心语义图谱主要关系统计}
\label{tab:kg_relation_stats}
\begin{tabular}{lrrrrrrrrr}
\toprule
关系 & 归属 & 导致 & 连接 & 发生于 & 关联监测 & 流向 & 处置 & 影响 & 表现为 \\
\midrule
数量 & 36 & 24 & 15 & 10 & 8 & 7 & 6 & 4 & 4 \\
\bottomrule
\end{tabular}
\end{table}

代表性知识链包括“HP级间汽封缺陷--蒸汽泄漏--汽耗率升高”“凝汽器水侧污垢--换热能力下降--凝汽器背压变化--LP有效焓降下降”“高压调节阀死区--阀位对指令不敏感--流量/压力振荡”等。这些关系均能够回溯到系统拓扑或原始研究证据\cite{ZhangHu2021FaultPropagationGNN,Medica2018CondenserCondition,ZhangHu2021ControlValveGraph}。当前图谱规模仅表示后续KG-GCN方法开发所需的\textbf{核心语义知识层}，随着本机组运行规程、缺陷记录、检修资料和全系统测点继续补充，节点和关系仍将持续扩展，不将107节点和114三元组视为最终工程知识库规模。

\begin{figure}[htbp]
\centering
\fbox{\parbox[c][5.0cm][c]{0.90\textwidth}{\centering
图3-3占位：汽轮机系统核心故障知识图谱\\
建议最终按主通流设备分区，突出调节阀死区、汽封泄漏、通流侵蚀/沉积、凝汽器污垢、给水加热器内漏等代表性故障链；\\
完整107节点/114边不全部标字，避免图面过密。}}
\caption{汽轮机系统核心故障知识图谱示意图}
\label{fig:turbine_core_kg}
\end{figure}

\subsubsection{领域语义层与DCS测量层映射}

领域知识图谱中的“VHP排汽温度”“一级抽汽压力”“凝汽器背压”等节点是稳定的领域参数概念，而实际DCS中可能存在A/B侧、多冗余传感器和具体KKS点号。若直接将所有点号作为唯一语义层，图谱将被大量设备专有命名和冗余测点占据，降低知识可读性和跨机组复用性。为此，本文采用“领域语义层+测量层”双层结构。

设领域语义图为$\mathcal{G}_{sem}$，参数实体$p_i\in\mathcal{V}_P$对应的实际测点集合表示为
\begin{equation}
\mathcal{M}(p_i)=\{m_{i1},m_{i2},\ldots,m_{in_i}\},
\label{eq:measurement_mapping}
\end{equation}
其中$\mathcal M$为领域参数到DCS测点的映射。语义层负责保存稳定的设备、故障和机理关系，测量层负责适配本文机组具体的传感器配置。对于多个冗余测点，可将其分别保留为测量节点并共同指向同一领域参数概念，也可在第四章根据KG-GCN输入节点定义选择最终测点级拓扑。

以第二章VHP为例，GRU实际预测27个候选状态，后续评价与诊断采用筛选后的24个状态点。这24个点首先映射至VHP本体状态、一级抽汽接口、VHP--RH1冷再接口以及缸体/内部蒸汽等领域参数概念，再通过“参数--归属--设备”“设备--介质流向--设备”以及“故障--关联监测--参数”等关系连接到全系统故障知识图谱。对于HP、IP、LP、再热和凝汽器等对象，后续同样可采用该方式将真实测点挂接到统一语义层。

由于本文故障诊断实验首先以VHP及其关联热力接口作为代表性对象，下一章可从全系统$\mathcal G_{sem}$中截取与VHP最终24个状态节点有关的参数子图，而不需要将整张汽轮机系统图的全部节点同时输入一个分类模型。这样既保持知识图谱的系统级完整性，又能够使当前故障实验与已有VHP状态监测结果直接衔接。

\begin{figure}[htbp]
\centering
\fbox{\parbox[c][4.2cm][c]{0.90\textwidth}{\centering
图3-4占位：上层为设备/参数/异常/故障/处置领域语义图 $\mathcal G_{sem}$；\\
下层为实际KKS/DCS测点；通过映射 $\mathcal M$ 将冗余测点挂接至领域参数；\\
VHP最终24状态点从测量层形成下一章KG-GCN参数子图。}}
\caption{领域语义层与DCS测量层映射关系}
\label{fig:semantic_measurement_mapping}
\end{figure}

\subsection{本章小结}

本章针对汽轮机故障知识分散于系统拓扑、DCS参数、运行资料和故障文献、难以直接被数据驱动模型利用的问题，构建了面向KG-GCN的汽轮机系统故障知识图谱方法。首先从主通流、再热、抽汽回热、给水和冷端系统分析故障传播路径，并结合原始文献整理调节阀死区、汽封泄漏、通流侵蚀/沉积、凝汽器污垢和给水加热器内漏等典型故障知识；随后定义设备、参数、异常、故障和处置五类实体，以及归属、介质流向、发生于、导致、影响、表现为、关联监测和处置等关系类型。

针对知识来源和构建阶段差异，本章将当前已完成的核心语义图谱与后续自动知识扩展流程进行区分。当前核心语义层以本机组系统拓扑、DCS参数概念和原始同行评议文献为证据，通过人工核验和schema规则化形成107个语义节点和114条关系三元组；对于后续运行日志、缺陷记录和检修文本，则设计大语言模型候选抽取、领域规则校验、Embedding语义对齐和多源证据评价流程进行扩展。跨机型文献中的知识仅作为通用热力机理来源，需通过本文机组结构和测点兼容性审核后使用。

进一步地，本章采用领域语义层$\mathcal G_{sem}$与DCS测量映射$\mathcal M$的双层设计，将稳定的汽轮机设备/故障知识与具体机组KKS测点分离。下一章将在该知识图谱基础上提取与状态监测节点相对应的参数图拓扑，并将实际运行状态与健康模型输出之间的动态偏差信息映射至图节点，形成基于KG-GCN的燃煤机组汽轮机系统故障诊断模型。
